\subsection{正则代数}

定义 1. --- 设 K 为一个域。称 K 上的代数 A 是正则的，如果对于 K 的每一个扩域 L，L ⊗K A 都是整环。

正则代数尤其是一个整环，因此是交换的。

命题 2. --- 设 A 和 B 是域 K 上的两个代数。若 A 是整环且 B 是正则的，则 \(A \otimes_{K} B\) 是整环。

设 \(L\) 为 \(A\) 的分式域。由于 \(B\) 是一个正则 \(K\) -代数， \(L \otimes_K B\) 是整环，因此 \(A \otimes_K B\) 亦然，因为它同构于 \(L \otimes_K B\) 的一个子环。

命题 3. --- 设 K 是一个域。

\begin{enumerate}
\def\labelenumi{\alph{enumi})}
\item
  正则 K-代数的每个子代数本身都是正则的。
\item
  两个正则 K-代数的张量积仍然是正则的。
\item
  设 A 是一个 K-代数，且 \(K'\) 是 K 的一个扩张。要使 A 为正则，必要且充分条件是由 A 通过标量扩张得到的 \(K'\) -代数 \(\mathrm{A}_{(\mathbb{K}')}\) 是正则的。
\end{enumerate}

\(a\) （相应地 \(c\) ) 的证明与 V，第 119 页命题 3 的部分 \(a\) （相应地 \(d\) ) 的证明相同，只需在处处将“既约环”替换为“整环”、“可分代数”替换为“正则代数”。我们来证明 \(b\) ).

设 A 和 B 是两个正则 K-代数。设 L 是 K 的一个扩张。由于 A 是正则的，环 \(L \otimes_{K} A\) 是整环。根据命题 2，环 \((L \otimes_{K} A) \otimes_{K} B\) 是整环，因为 B 是正则的。最后，环 \(L \otimes_{K} (A \otimes_{K} B)\) 同构于 \((L \otimes_{K} A) \otimes_{K} B\) ，因此是一个整环。这表明 \(A \otimes_{K} B\) 是一个正则 K-代数。

